% Draft manuscript outline for Journal of Chemical Information and Modeling.
% This is intentionally an outline and is NOT ready for journal submission.
% Scope: tetrahedral stereochemistry; the input graph already contains every
% candidate stereodescriptor (parity=None means unassigned).
\documentclass[journal=jcisd8,manuscript=article]{achemso}
% Corresponding-author email is deliberately omitted in this preliminary outline.
% Before submission, remove email=false and supply a validated \email{...}.
\setkeys{acs}{email=false}

\usepackage{amsmath,amssymb}
\usepackage{booktabs}
\usepackage{graphicx}
\usepackage{array}
\emergencystretch=2em % improve line breaks in this long outline

\newcommand{\StereoMolGraph}{\textsc{StereoMolGraph}}
\newcommand{\Aut}{\operatorname{Aut}}
\newcommand{\Stab}{\operatorname{Stab}}
\newcommand{\sgn}{\operatorname{sgn}}
\newcommand{\Iso}{\cong}
\newcommand{\todolist}[1]{\par\smallskip{\small\noindent\textbf{To develop:} #1\par}\smallskip}

\title{Conditional Stereogenicity and Stereoisomer Enumeration Using Molecular Graph Isomorphisms}

% Author list for the working draft; affiliation to be added later.
\author{Maxim Papusha}

\begin{document}

\begin{abstract}
Some tetrahedral atoms become stereogenic only after configurations at other
sites are specified. This dependence is difficult to capture by examining
individual substituents. We present a graph-based framework that compares
alternative local configurations while holding the remaining stereochemistry
fixed. It identifies always-active, conditionally active, and inactive
tetrahedral sites and records which remote assignments change their
stereogenicity. Complete configurations are generated from existing
stereodescriptors and compared by stereo-aware graph isomorphism. In an
illustrative case, a symmetric oxygen-bridged molecule yields seven distinct
stereochemical graph classes, confirmed by an independent automorphism
calculation, whereas a standard RDKit enumeration yields six. The approach
provides a basis for analyzing stereochemical dependence and testing
stereoisomer enumeration, subject to the chemical information and
geometric constraints represented by the input graph.

\end{abstract}

% TODO: Final ACS submission needs a table-of-contents graphic (TOC entry),
% final author metadata, a polished abstract, and separate SI files.

\section{Introduction}
\label{sec:intro}

A tetrahedral carbon attached to four different substituents is a familiar
example of a stereogenic center. However, equivalent substituents can
become distinguishable after configurations elsewhere in the molecule
are specified. This is particularly relevant in symmetric and polycyclic
structures. Examining one atom without its stereochemical surroundings
may therefore give an incomplete answer. IUPAC terminology already
distinguishes a stereogenic unit from a chiral \emph{molecule}.\cite{Moss1996,FavrePowell2014}

Consider 2,3-dichlorobutane. Its two tetrahedral centers permit four
formal combinations of local configurations. Molecular symmetry makes
$(2R,3S)$ and $(2S,3R)$ descriptions of the same meso isomer, so only
three stereoisomers remain. This separates two questions:
\emph{Does changing one selected center matter when the others are fixed?}
and \emph{Are two fully configured molecules equivalent?}
Cahn--Ingold--Prelog (CIP) labels name configurations but do not
themselves answer these graph-equivalence questions.\cite{FavrePowell2014,Hanson2018}

Established software includes RDKit's perception of potential stereo
sites,\cite{RDKitBook2026} graph-automorphism-based stereoisomer
enumeration in \textit{enu},\cite{Rieder2023} and
stereopermutation-based representations such as
\textit{Molassembler}.\cite{Sobez2020}
The previously published \StereoMolGraph{} library incorporates
local permutation-group descriptors into molecular and reaction
graph comparisons.\cite{Papusha2026} This study does not
reintroduce that representation, assign CIP priorities, or infer
coordination geometries.

Instead, we investigate how the \emph{distinguishability} of a
tetrahedral center changes with configurations at other sites.
We define a site-specific graph comparison, use it to build
stereochemical dependency networks, and test exhaustive
stereoisomer enumeration against independent symmetry calculations.
All potential tetrahedral sites are supplied with the input:
the algorithms never add new descriptors.

\section{Theory and Mathematical Definitions}
\label{sec:theory}

\subsection{Input molecular graph and local configuration domains}
\label{subsec:input}

The input is a molecular graph containing a list of potential
tetrahedral sites. We describe it by
\begin{equation}
G=(V,E,\lambda_V,\lambda_E,\mathcal D),
\label{eq:graph}
\end{equation}
where $V$ is the set of atoms, $E$ is the set of bonds,
$\lambda_V$ and $\lambda_E$ are the stored chemical attributes,
and $\mathcal D=\{d_s:s\in S\}$ contains the pre-existing
tetrahedral descriptors. In chemical terms, the input already
specifies which atoms have tetrahedral geometries that may be
varied. The upstream adapter supplies these sites; detection
and enumeration do not introduce any new descriptors.

A tetrahedral descriptor has two formal spatial arrangements,
denoted $+1$ and $-1$ relative to its ordered ligand positions.
These are local parity labels, \emph{not} the CIP descriptors
$R$ and $S$.\cite{Hanson2018,Papusha2026}
Let $X_s=\{-1,+1\}$ contain both local states at an existing
site. If one state is already specified, it is held fixed during
enumeration. Thus the allowed choices are
\begin{equation}
D_s =
\begin{cases}
\{\text{specified state}\}, & \text{if $s$ is fixed},\\
X_s, & \text{if $s$ is unspecified}.
\end{cases}
\label{eq:domain}
\end{equation}
The complete set of formal assignments is
\begin{equation}
\Omega=\prod_{s\in S}D_s,\qquad
|\Omega|=\prod_{s\in S}|D_s|.
\label{eq:omega}
\end{equation}
For five unspecified tetrahedral sites, this gives $2^5=32$
assignments. Some assignments may represent identical
stereochemical structures because of symmetry. Others may
be incompatible with ring geometry. Equation~\ref{eq:omega}
therefore describes the search space, not necessarily the
number of chemically distinct stereoisomers.

\subsection{Site-anchored stereogenicity in a fixed context}
\label{subsec:local-test}

Consider one potential tetrahedral center $s$. Fix all
other stereo configurations and ask whether exchanging the
two local arrangements at $s$ changes the stereochemical
structure. Two substituent branches that appear identical
before stereo is specified may become distinguishable after
their remote configurations are assigned.

We answer this question using \emph{stereo-preserving graph
isomorphism}: a one-to-one correspondence between atoms
that preserves their connectivity, stored chemical attributes,
and local stereochemical arrangements.\cite{Papusha2026,Cordella2004}
For a \emph{site-specific} comparison, the chosen atom $s$
must also map onto itself. We denote this by $\cong_s$.
This condition prevents an otherwise valid global
isomorphism from exchanging $s$ with a different,
symmetry-equivalent atom.

Let $\boldsymbol{\sigma}_{-s}$ denote the fixed
configurations at all other sites.
The graph $G_{s,x}(\boldsymbol{\sigma}_{-s})$ has
local state $x$ at $s$. Define
\begin{equation}
A_s(\boldsymbol{\sigma}_{-s}) =
\begin{cases}
1, & G_{s,+1}(\boldsymbol{\sigma}_{-s})
       \not\cong_s G_{s,-1}(\boldsymbol{\sigma}_{-s}),\\
0, & G_{s,+1}(\boldsymbol{\sigma}_{-s})
       \cong_s G_{s,-1}(\boldsymbol{\sigma}_{-s}).
\end{cases}
\label{eq:activation}
\end{equation}
Here $A_s=1$ means that the two arrangements cannot
be exchanged without changing the represented stereo graph
\emph{while keeping the chosen atom fixed}.
To assess this property one may compare both hypothetical
states of $s$ even if the input currently fixes one state;
the configurations at other sites remain fixed.

Repeat the test for all allowed configurations of the
other sites. This gives three classes:
\begin{equation}
\begin{aligned}
\text{always active:}\quad & A_s(\tau)=1\text{ for every }\tau,\\
\text{conditional:}\quad & A_s(\tau)=0\text{ for some }\tau
                         \text{ and }1\text{ for others},\\
\text{inactive:}\quad & A_s(\tau)=0\text{ for every }\tau .
\end{aligned}
\label{eq:class}
\end{equation}
The symbol $\tau$ denotes one complete choice at the
other supplied sites. These classifications concern the
\emph{represented graph}, not physical stability or
automatically assigned CIP descriptors. The anchored
local question must also be distinguished from equivalence
of two complete molecules without any anchored atom.

\subsection{Directed stereochemical dependence}
\label{subsec:dependency}

A configuration at site $t$ \emph{influences} another site $s$
if changing $t$ can switch the activity of $s$ while all
other configurations remain unchanged. We represent this by
a directed arrow $t\rightarrow s$. The opposite arrow
need not exist.

Such an arrow exists if there is at least one assignment of
the remaining sites for which
\begin{equation}
\begin{aligned}
t\longrightarrow s\quad\Longleftrightarrow\quad
&\exists\,\boldsymbol{\tau}_{-\{s,t\}},\ x,y\in X_t:\\
&A_s(\boldsymbol{\tau}_{-\{s,t\}},t=x)
\ne A_s(\boldsymbol{\tau}_{-\{s,t\}},t=y).
\end{aligned}
\label{eq:dependency}
\end{equation}
The symbol $\exists$ means that at least one such case exists.
Each arrow can be supported by an explicit pair of
assignments, called a \emph{witness}. The arrow does
not mean that every change at $t$ affects $s$.

The sites and arrows form a directed dependency
graph $\mathcal H=(S,E_{\mathcal H})$. In the
cyclic molecule VS294, the preliminary analysis
finds four arrows into C3, from C1, C5, C8, and C9.
Changing any one source can change the activity
of C3 in some fixed stereochemical context.

A dependency graph tells us \emph{which} sites
may matter. The full function $A_s$ also tells
us \emph{when} they matter, including joint
effects of multiple configurations. Molecular
color refinement offers a fast way to suggest
possible influences,\cite{Morgan1965} but equal
colors do not prove ligand equivalence.
Heuristic edges must not be used to discard
stereo assignments.

\subsection{Atom orbits, chirality, and stereochemical context}
\label{subsec:orbits}

Two atoms belong to the same \emph{orbit} if a symmetry
of the molecular graph can exchange them. For example,
the two middle carbon atoms of succinic acid are
constitutionally equivalent. A symmetry of a graph is
called an \emph{automorphism}: it exchanges equivalent
atoms while retaining their bonds and chemical labels.

After tetrahedral configurations are specified, only
automorphisms preserving those configurations remain
allowed. The atom orbit is
\begin{equation}
\mathcal O_{\sigma}(v)
=\{g(v):g\in\Aut(G;\sigma)\},
\label{eq:orbits}
\end{equation}
where $\Aut(G;\sigma)$ is the group of stereo-preserving
automorphisms. Assigning stereochemistry can split
previously equivalent atom orbits but cannot merge
previously distinct ones when the underlying labels
stay unchanged.

\emph{Atom equivalence is not the same as molecular
chirality.} A molecule is chiral when its mirror image
is not superposable on the original.\cite{Moss1996,FavrePowell2014}
A meso molecule may contain stereogenic centers yet be
achiral. Conversely, equivalent atom positions can remain
related by symmetry even when the completed molecule is
chiral. To test molecular chirality one compares a complete
assignment with its inverted, mirror assignment, without
anchoring any one site.

\todolist{Add independently verified chemical examples
of an atom orbit being preserved or split by stereochemical
decoration. Calculate their automorphism groups before
reporting the orbit sizes.}

\subsection{Enumerating unique stereochemical graph classes}
\label{subsec:enumeration}

The reference generator visits every assignment in
eq~\ref{eq:omega}, sets the configurations at the
\emph{existing} descriptors, and removes equivalent
complete stereo graphs by graph isomorphism. Unlike
an initial local ligand test, this comparison accounts
for all other assigned configurations.

Molecular symmetry provides an independent way to
check the count. A constitutional automorphism may
exchange equivalent stereo sites. It also transports
their local tetrahedral configurations, possibly
reversing a parity sign when the ligand ordering is
changed. All complete assignments connected by these
operations form a \emph{symmetry orbit}, representing
one stereo-graph equivalence class.

If $\Gamma$ is the group of structural symmetries
compatible with fixed input stereochemistry, Burnside's
formula gives the number of distinct graph classes:
\begin{equation}
N_{\mathrm{orbits}}=
\frac{1}{|\Gamma|}
\sum_{g\in\Gamma}|\operatorname{Fix}_{\Omega}(g)|.
\label{eq:burnside}
\end{equation}
To apply it, count the assignments unchanged by
each symmetry, add those counts, and divide by the
number of symmetries.\cite{Rieder2023}
For 2,3-dichlorobutane, the two equivalent meso
descriptions occupy one orbit, leaving three
stereoisomers instead of four assignments.

\subsection{Prospective symmetry-pruned enumeration}
\label{subsec:prune}

The baseline removes duplicates only after full
assignments are constructed. A future algorithm
could avoid exploring equivalent intermediate
assignments using canonical augmentation or
stabilizer-orbit techniques.\cite{McKay1998}
For a partial assignment $\sigma_j$ and its
next site $s_j$, consider the symmetries
that preserve both:
\begin{equation}
H_j=\Stab_{\Gamma}(\sigma_j)\cap\Stab_{\Gamma}(s_j).
\label{eq:stabilizer}
\end{equation}
Here $\Stab$ denotes a subgroup whose elements
leave the specified object unchanged. Under an
exact and validated symmetry action, one may
explore representatives of the local-state
orbits instead of all equivalent choices.
This acceleration is not implemented in the
reference generator. Any proposed pruning must
reproduce its complete equivalence classes.

\section{Computational Methods and Reproducibility}
\label{sec:methods}

\subsection{Input preparation and scope}

All calculations use \StereoMolGraph{} instances with
pre-existing potential tetrahedral descriptors.
An unspecified configuration is represented by
\texttt{parity=None}; a specified configuration stays
fixed. Analysis and enumeration do not add stereo
sites. For the initial comparison, molecules with
isotopically labeled atoms or non-single bonds
(including aromatic bonds) are excluded, because
these labels are not yet fully represented in the
graph attributes used for matching. This is an
explicit restriction of the benchmark, not a
general restriction of the theory.

Source SMILES in the pilot are drawn from the
300-molecule CIP validation suite published with
the algorithmic analysis of the CIP rules.\cite{Hanson2018}
That suite tests nomenclature and descriptor
assignment and is not an independent reference
set of stereoisomer counts. For a controlled
comparison, the input adapter retains only
\emph{existing} tetrahedral descriptors that
match RDKit's potential stereo-site list.
This isolates enumeration from upstream
site nomination. RDKit's stereoperception
and enumeration are used as comparators,
not as an assumed ground truth.\cite{RDKitBook2026,RDKit2026}

\subsection{Local activity and dependency tests}

For each site, we hold the configurations
at all other sites fixed, compare the two
states of the selected tetrahedron, and
apply the anchored stereo-aware graph
isomorphism test of eq~\ref{eq:activation}.
The implementation uses the stereochemical
graph comparison in \StereoMolGraph{} and
a VF2-family search.\cite{Papusha2026,Cordella2004}
Identical contexts are cached so that a
site is not retested unnecessarily.

To determine an exact dependency arrow
$t\rightarrow s$, we search for a pair
of otherwise identical complete assignments
that differ at $t$ and give different values
of $A_s$. Each edge is stored with its
witness assignments. A faster graph-color
test can suggest positional influences,
but no approximate edge is used to remove
a configuration.\cite{Morgan1965}

\subsection{Independent counting of stereochemical classes}

For the small pilot cases, construct an
explicit-hydrogen graph and calculate its
constitutional automorphisms using
NetworkX.\cite{Hagberg2008}
Transport the ligand ordering and tetrahedral
parity under each automorphism, then
partition all local assignments into orbits.
The orbit count from eq~\ref{eq:burnside}
is checked against direct pairwise
stereo-aware isomorphism of the generated
representatives. Thus, a count matching
RDKit is not sufficient to establish
correctness: the generated equivalence
classes must agree.

\subsection{Geometry and stereochemical cleanup}

For any discrepancy with RDKit, compare
tetrahedral tags before and after RDKit's
stereochemistry cleanup (\texttt{cleanIt=True}).
A preliminary geometric check uses RDKit's
ETKDGv3 conformer generation, including
the developments for cyclic molecules,
followed by MMFF minimization.\cite{Riniker2015,Wang2020,Halgren1996}
Successful embedding supports the
existence of a locally consistent 3D
arrangement, but does not establish
configurational stability, a reaction
barrier, or experimental isolability.

\subsection{Planned validation and scaling}

\todolist{Assemble a curated tetrahedral-only
benchmark that includes acyclic molecules,
meso compounds, repeated stereogenic branches,
bridged rings, cages, and conditional centers.
Verify each complete equivalence class using
an independent symmetry calculation.
Test invariance under atom relabeling
and consistent reordering of ligand slots.
Report runtime, graph-isomorphism calls,
visited search nodes, and peak memory
for the exhaustive baseline and any
future symmetry-pruning implementation.}

\section{Results and Discussion: Working Outline}
\label{sec:results}

\subsection{A case study with five tetrahedral sites}

The oxygen-bridged molecule VS294 is part of the published
CIP validation suite.\cite{Hanson2018} It has five supplied tetrahedral
sites, numbered 1, 3, 5, 8, and 9 in the input graph. For this
molecule, sites 1, 5, 8, and 9 are always active under
eq~\ref{eq:activation}, whereas C3 is conditionally active.
In particular, changing the configuration at any one of
C1, C5, C8, or C9 changes the activity of C3 for at
least one otherwise fixed assignment. These are
\emph{configuration-dependent} relationships, not simply
the locations of neighboring atoms.
Figure~\ref{fig:dependency} will show the molecular structure
and the four verified directed arrows.

\begin{figure}[ht]
  \centering
  \fbox{\parbox{0.87\linewidth}{\centering\small
  \vspace{0.6cm}
  PLACEHOLDER: VS294 molecular drawing with atom indices and verified
  directed dependencies toward central site C3.
  \vspace{0.6cm}}}
  \caption{Planned illustration of conditional tetrahedral stereogenicity
  in VS294. Solid edges will represent exact, witness-supported changes
  in the stereogenicity of the central site, rather than heuristic
  positional influence.}
  \label{fig:dependency}
\end{figure}

\subsection{Seven graph classes versus six RDKit outputs}

The same five variable sites give 32 local tetrahedral assignments.
An independent explicit-hydrogen graph automorphism calculation yields
$|\Gamma|=8$ and seven disjoint assignment orbits with sizes
$2,2,4,4,4,8,8$. The Burnside fixed-point sum is 56, hence seven
classes. All seven \StereoMolGraph{} representatives occupy different
independently calculated orbits and all 21 pairs are non-isomorphic
under the current stereo-aware VF2 implementation.

In the pilot, RDKit 2026.09.2 generates six unique outputs.
Its standard stereochemical cleanup removes the C3 tetrahedral
tag in all 32 explicitly tagged assignments, merging
two otherwise distinct orbits. The independent group
calculation identifies seven classes: three enantiomeric
pairs and one achiral class in the represented stereo model.
Ten conformers of each of two opposite C3 configurations
were embedded and optimized by molecular mechanics; the
other four local configurations remained unchanged.
This supports local 3D realizability but not the
experimental isolability or interconversion barriers
of these structures.\cite{RDKit2026,Riniker2015,Wang2020,Halgren1996}

\begin{table}[ht]
\centering
\caption{Verified pilot results for VS294 under the stated
 tetrahedral graph model; not a full benchmark.}
\label{tab:vs294}
\begin{tabular}{@{}lr@{}}
\toprule
Quantity & Result \\
\midrule
Candidate tetrahedral sites & 5 \\
Formal assignments & 32 \\
Constitutional automorphisms & 8 \\
Independent stereo-graph orbits & 7 \\
StereoMolGraph representatives & 7 \\
Distinct RDKit-cleaned outputs & 6 \\
Pairwise isomorphic pairs among seven & 0/21 \\
\bottomrule
\end{tabular}
\end{table}

\subsection{Accuracy and generalization benchmark}

\todolist{Replace pilot table with a well-controlled benchmark and
explicit characterization of which molecules are in scope. Preliminary
isotope-free, single-bond comparisons matched RDKit counts in 137 of
152 completed cases, with 15 mismatches concentrated in cyclic
structures; this is a selected, RDKit-site-matched diagnostic sample,
not a general accuracy estimate. Supply an independent equivalence
oracle for each nontrivial discrepancy.}

\subsection{Symmetry, meso compounds, and activation patterns}

\todolist{Quantify always-active, conditionally active, and inactive
sites by topology. Classify enantiomeric pairs and achiral
stereo-graph classes. Test atom-orbit preservation/splitting under
stereochemical decoration and distinguish it from molecular chirality.}

\subsection{Computational scaling and symmetry pruning}

\todolist{Benchmark the exhaustive reference implementation against
incremental caching and certified automorphism pruning. Report results
as a function of both the number of candidate sites and the number of
inequivalent outputs. Do not claim acceleration until correctness is
validated against exhaustive enumeration on the same inputs.}

\section{Limitations and Scope}

The current results concern tetrahedral configuration domains
associated with \emph{given} molecular graph descriptors. Candidate
nomination and chemically valid valence/hybridization perception are
separate problems. Exact graph isomorphism establishes equivalence
relative to the represented atom and bond attributes; missing
constitutional labels can merge distinct chemical species. The
existence of a geometrically embedded conformer does not prove
configurational stability or experimental isolability. Isotopic,
alkene, axial, and coordination stereochemistry remain outside the
initial benchmark scope.

\section{Conclusions: Intended Claims}

A rigorous description of configuration-dependent stereogenicity
allows molecular sites to be classified by how their distinguishable
states change under assignments elsewhere in the molecule. Exact
anchored isomorphism tests turn such effects into reproducible
dependency witnesses, and the resulting configured graphs can be
validated independently through automorphism orbits. The full
manuscript must establish generality and computational effectiveness
through a broader benchmark before drawing conclusions beyond the
current case studies.

\section{Data and Software Availability}

Source code and example analyses are available in the GitHub projects
\emph{StereoMolGraph} and \emph{Experiments-Stereoisomers}
(owner: \texttt{maxim-papusha}).
\todolist{Before submission: pin immutable release commits, deposit
full benchmark tables and reproducibility scripts, supply software
license and environment information, and archive a release DOI.}

\section{Supporting Information}

\todolist{Prepare a separate Supporting Information PDF or document
containing: formal proofs and descriptor transport conventions;
complete assignment/orbit tables; atom-numbered structures;
all graph-automorphism witnesses; benchmark membership and exclusion
criteria; exhaustive vs optimized enumeration verification;
relabeling invariance checks; and conformer-generation details.}

% Confirm authors and funding acknowledgments before submission.
\bibliography{references}

\end{document}